\section{Primes in a prescribed ray class on short intervals}
\label{sec:ray-prime-estimates}

The auxiliary level requires a rational prime in a short interval with a
prescribed representation by a quadratic norm. Fix a prime $d\ge100$ with
$d\equiv3\pmod4$, let $K=\mathbb Q(\sqrt{-d})$, and fix an integer
$M\ge3$. A prime ideal of norm $p$ in the identity ray class modulo
$M\mathcal O_K$ has a generator $\pi\equiv1\pmod{M\mathcal O_K}$.
Writing out $N_{K/\mathbb Q}(\pi)=p$ gives integers $A,D$ with
\[
 4p=A^2+dD^2,\qquad M\mid D,\qquad 2M\mid A-D-2 .
\]
The modulus varies with the level. We need a lower bound for the number
of such primes which still says something as the modulus grows.
Put $\mathcal B=4dM^2$. We shall obtain absolute constants $c,C>0$ such that
the number of these primes in $(X,X(1+h)]$ is at least
\[
   C\,\frac{hX}{\mathcal B^{47}\log X}
   \quad\text{if}\quad
   \log\mathcal B\le c\sqrt{\log X},
   \qquad e^{-c\sqrt{\log X}}\le h\le1 .
\]

There are two points in the proof where uniformity matters. Counting ideals
in a fixed ray class amounts to counting points of a translated lattice in an
ellipse. The error must be uniform in the translation. Later, orthogonality
expresses the prime sum in the identity class as an average over ray-class
characters. A real zero near $1$ could use up almost all the main term for a short
interval. We show that at most one such zero occurs in the whole family and
bound its distance from $1$ before summing the errors over the characters.
Removing higher prime powers then gives the count of rational primes.

In what follows $\mathcal B$ is always the conductor bound. The earlier scale
$B=e^{x^\epsilon}$ has a different meaning. We use multiplicative notation for
ideals and sometimes additive notation for the ray class group, writing $0$ for
its identity in sums over characters.


\subsection{Lattice points and the ray class group}

\paragraph{Ideals as points of a lattice coset (QR1).}

To choose the representative below, take an integral ideal $\mathfrak a$ in the inverse class, coprime to $M\mathcal O_K$. Then $\mathfrak a+M\mathcal O_K=\mathcal O_K$ gives $\alpha_0\in\mathfrak a$ congruent to $1$ modulo $M\mathcal O_K$. Since $\mathfrak a\cap M\mathcal O_K=M\mathfrak a$, every such element lies in the same coset $\alpha_0+M\mathfrak a$.

% excerpt: Let $d\ge100$ be prime with $d\equiv3\pmod4$, let $M\ge3$, and set || no unit multiplicity to absorb in a later leading constant.


\paragraph{Reduction to a binary quadratic form (QR2).}

% excerpt: Put $N=N\mathfrak a$ and $\Delta=\operatorname{covol}(M\mathfrak a)$. || ideals of norm $j$, prime to $M$, in the original class $g$.


\paragraph{The count in an ellipse (QR3).}

% excerpt: For $x\ge1$ and fixed $n$, the condition $Q(m,n)\le x$ is an interval || the two shifts changes this estimate.


\paragraph{A Gaussian sum with a shift (QR4).}

% excerpt: For $s>0$, periodize $e^{-s t^2}$ on $\mathbb R/\mathbb Z$. || can be used inside a second Gaussian sum.


\paragraph{The ordinary and Gaussian ideal counts (QR5).}

% excerpt: Apply QR4 to the inner sum in $m$ in $e^{-tQ(m,n)}$; its shift || modulus.


\paragraph{The convolution from ideal factorization (QR6).}

% excerpt: Define $\lambda(\mathfrak p^r)=\log N\mathfrak p$ for a nonzero || convergent series occurs here.


\paragraph{Prime norms and the congruence condition (QR7).}

% excerpt: All the weights are nonnegative. A nonzero weight at an integer norm || replacement of a ray condition by ordinary principalness has been made.


\paragraph{Positivity for real characters (QR8).}

% excerpt: Let $\chi^2=1$ be a ray character, so $\chi$ takes values $\pm1$. || positivity of a few initial coefficients.


\paragraph{The order of the ray class group and the family bounds (QR9).}

% excerpt: The ideal of norm one is $\mathcal O_K$, in the identity class; || as a hypothesis of QW11.


\subsection{The associated $L$-functions near one}

\paragraph{Continuation of the character sums (QR10).}

% excerpt: For the rest of this section let $\mathcal F$ be any \texttt{LFam}, || $F_\chi(\bar s)=\overline{F_\chi(s)}$ by conjugating the integral.


\paragraph{Nonvanishing to the right of one (QR11).}

% excerpt: The nonnegative weights and their total bound give || presupposing an Euler product for the abstract family.


\paragraph{The heat sum at one (QR12).}

% excerpt: Put $\Theta_\chi(t)=\sum_n a_\chi(n)e^{-tn}$ for $t>0$. ||                           =m_\chi(1)=F_\chi(1).           \tag{QR.25}
\]

\paragraph{A lower bound at one for a real character (QR13).}

% excerpt: Let $\chi\ne1$ be real. Write $P=\Theta_1$, $Q=\Theta_\chi$, || this step does not invoke an ineffective exceptional-zero theorem.


\paragraph{Separation of a real zero from one (QR14).}

% excerpt: If $|z-1|\le1/4$, the disk $|w-z|\le1/10$ lies in || conductors.


\subsection{A zero-free region uniform in the family}

\paragraph{The local logarithmic derivative (QR15).}

% excerpt: Here is the complex-analysis estimate used in the zero-free argument. || no unproved global zero-count estimate is substituted for it.


\paragraph{Uniform bounds for the local zero sum (QR16).}

% excerpt: Put $\ell(t)=\log(\mathcal B(|t|+2))>2$ and center a disk at || $\ell(2t)\le2\ell(t)$ will be used below.


\paragraph{The zero-free region (QR17).}

% excerpt: For $|z|=1$, || that the next steps control.


\paragraph{Uniqueness of the exceptional zero (QR18).}

% excerpt: Increase the constant $C$ of (QR.34) if needed, and choose || not have supplied that assertion.


\paragraph{The logarithmic derivative after removing the exception (QR19).}

% excerpt: Fix $0<c_0\le1/20$ and a finite set $Z$ satisfying the literal || also proves $F_\chi(s)\ne0$ at the points under consideration.


\subsection{Counting primes on the short interval}

\paragraph{Truncated Perron inversion (QR20).}

% excerpt: Let $b(n)$ be any complex coefficients with $|b(n)|\le2\Lambda(n)$. || floor convention, without a hidden distance-to-integers restriction.


\paragraph{The exceptional residue (QR21).}

% excerpt: Assume $\operatorname{ExcZ}(\chi,c_0,Z)$ as in QR19, let $u\ge2$ || multiply that main term.


\paragraph{The height of the integration contour (QR22).}

% excerpt: Suppose $\log\mathcal B\le\sqrt{\log u}$ and put || range; the source's sharper value is not needed for its conclusion.


\paragraph{One exceptional term for the family (QR23).}

% excerpt: Choose once || this argument.


\paragraph{The identity-class sum on a short interval (QR24).}

% excerpt: Take $A_0=46$ and the constants $c_9,c_7,C_7$ just proved. Choose || short intervals.


\paragraph{The resulting prime count (QR25).}

% excerpt: Insert QR9's actual family, with $\mathcal B=4dM^2$ and norm-one || objects required later in the RH and Dirichlet-GRH branches.
